\documentclass[12pt,a4paper]{article}

\usepackage[margin=2.5cm]{geometry}
\usepackage{amsmath,amssymb}
\usepackage{graphicx}
\usepackage{booktabs}
\usepackage{hyperref}
\usepackage[numbers,sort&compress]{natbib}
\usepackage{setspace}
\usepackage{microtype}
\usepackage{xcolor}
\usepackage{enumitem}
\usepackage{caption}
\usepackage{abstract}
\usepackage{titlesec}
\usepackage{array}
\usepackage{multirow}

\onehalfspacing
\hypersetup{colorlinks=true,linkcolor=blue!60!black,
            citecolor=blue!60!black,urlcolor=blue!60!black}
\titleformat{\section}{\large\bfseries}{\thesection.}{0.5em}{}
\titleformat{\subsection}{\normalsize\bfseries\itshape}{\thesubsection.}{0.5em}{}
\setlength{\parindent}{1.5em}
\setlength{\parskip}{0.3em}

\title{\textbf{Toward a Measurable Theory of Complexity:\\
Eight Candidate Laws, Metric Derivation,\\
and Empirical Validation Across Four Simulation Substrates}\\[0.8em]
\large\textit{A report inviting expert scrutiny and collaboration}}

\author{Cameron Belt\\
\small Belt Software Solutions\\
\small \href{mailto:}{[contact via correspondence]}}

\date{\today\\[0.3em]
\small\textit{Preprint --- not peer reviewed}}

\begin{document}
\maketitle

%% ─── Abstract ───────────────────────────────────────────────────────────────
\begin{abstract}
\noindent
We present a framework for measuring complexity derived from eight proposed
necessary properties of genuinely complex systems, and document how each
metric weight was discovered through iterative empirical testing rather than
assumed from theory.
The framework operationalises four of the eight laws as measurable metrics:
spatial opacity (P1), gzip compression ratio (P2), spatial entropy with an
attractor weight (P7), and temporal compression with a dual-Gaussian weight
(P4/P5), combined multiplicatively so that a system must score well on all
dimensions simultaneously.

The central methodological contribution is the design of
\emph{initial-condition-independent weight functions}.
The original entropy weight (a fixed Gaussian calibrated to single-cell IC)
was replaced with an attractor formulation using $\tanh$ gates plus a
variance bonus on entropy standard deviation $\sigma_H$, after discovering
that Class~4 and Class~3 rules are indistinguishable by mean entropy alone
under random initial conditions.
The temporal compression weight was made IC-independent via a dual-Gaussian
that rewards both known Class~4 attractor values (0.58 under random IC;
0.73 under single-cell IC) without parameter switching.
A principled finding about opacity is also reported: it is contaminated by
the triangular cone transient under single-cell IC and is valid only under
random initial conditions, which are adopted as the universal experimental
standard.

Applied to all 256 Wolfram elementary cellular automata ($W{=}150$,
random 50\% IC, 5 seeds), the composite achieves \textbf{45$\times$}
Class~4 / Class~3 separation with all four known Class~4 rules at ranks 1--4
of 256.
Scale invariance is confirmed across grid widths 100--300.
Extension to radius-2 totalistic 1D CA (32 rules) yields 21$\times$
separation with correct ordering.
Extension to 2D Life-like rules (17 rules) places Conway and HighLife in
the top 10, with a diagnosed and addressable calibration gap in temporal
compression for two-dimensional substrates.
An N-body parameter scan (256 force-constant combinations) finds a structured
complexity landscape with a sharp boundary consistent with analytical
predictions; our universe sits at the 54th percentile.

We report both successes and failures honestly and invite expert scrutiny.

\bigskip
\noindent\textbf{Keywords:} complexity measures, cellular automata,
edge of chaos, IC-independent metrics, attractor weight, entropy variance,
scale invariance, Life-like rules, N-body, fine-tuning
\end{abstract}

\rule{\linewidth}{0.4pt}

%% ─────────────────────────────────────────────────────────────────────────────
\section{Introduction and Motivation}
%% ─────────────────────────────────────────────────────────────────────────────

The question of how to measure complexity is not new, but remains genuinely
open.
Existing proposals --- algorithmic complexity \citep{kolmogorov1965},
thermodynamic depth \citep{lloyd1988}, integrated information
\citep{tononi2004}, energy rate density \citep{chaisson2001} --- each capture
something real while failing to capture everything.
No single metric produces rankings that align with intuitive complexity
judgements across system types \citep{gell-mann1996,ladyman2013}.

The work reported here began not as a formal research programme but as a
philosophical inquiry: \emph{can the properties that make a system genuinely
complex be stated precisely enough to derive measurable metrics from first
principles?}
Through iterative reasoning, eight candidate laws were developed, each
describing a property that appears necessary for genuine complexity.
These were operationalised as computational metrics and tested against known
ground truths in several dynamical system classes.

This paper is structured to reflect how the work actually proceeded ---
from philosophical motivation, through metric derivation and iterative
refinement, to the final validated results.
We believe this structure is more useful than a conventional
hypothesis-results format for a framework paper, because the refinement
process itself contains the key methodological contributions.
In particular, the discovery of the entropy variance signal and the diagnosis
of opacity IC-contamination both emerged from the refinement process and
are not findings that could have been anticipated from theory alone.

The paper proceeds as follows.
Section~2 states the eight candidate laws.
Section~3 derives the metrics from those laws and documents the full
iterative refinement process by which each weight function was arrived at.
Sections 4--7 report experiments across four simulation substrates.
Section~8 discusses what the results establish, what they do not, and
the most important open questions.

%% ─────────────────────────────────────────────────────────────────────────────
\section{Eight Candidate Laws of Complexity}
%% ─────────────────────────────────────────────────────────────────────────────

The following eight properties are proposed as necessary conditions for
genuine complexity.
They are stated descriptively --- relationships derived from observation and
reasoning, not from more fundamental principles.
No single property is proposed as sufficient; the claim is that all are
necessary.

\subsection{P1 --- Opaque Hierarchical Layering}

A system is complex of degree $N$ when its state space is partitioned
into $N$ stable hierarchical layers such that:
\begin{equation}
H(S_K \mid S_{K+1}) > \varepsilon \quad \text{for all } K \in \{1,\ldots,N-1\}
\end{equation}
where $H(\cdot|\cdot)$ is conditional entropy.
Each higher layer is a lossy compression of the layer below; lower-level
information is irreversibly hidden from higher-level description.

\subsection{P2 --- Modular Interconnection}

Complexity scales multiplicatively with semi-autonomous modular architecture
within each layer.
Modules must be internally coherent but loosely coupled, so their interactions
produce behaviour that no individual module contains.
Empirically, complex systems tend toward scale-free small-world topology
\citep{watts1998,barabasi1999}.

\subsection{P3 --- Self-Assembly Without Top-Down Specification}

Genuine complexity cannot be engineered top-down.
Any fully specified interface ceases to be genuinely opaque (violating P1).
The most complex known systems arose through processes that did not require
a designer to understand the result.

\subsection{P4 --- Recursive Self-Modification}

Maximum complexity is bounded by capacity to modify the system's own
generative rules.
Fixed-rule systems can explore but not expand their state space.
Only recursive self-modification can generate new abstraction layers.

\subsection{P5 --- Selection Pressure}

Self-assembling systems require asymmetric environmental constraints to
generate directed complexity rather than random drift.
The richness of complexity generated scales with the intensity, variety,
and responsiveness of selection pressures.

\subsection{P6 --- Co-evolutionary Dynamics}

Maximum complexity is determined not only by internal properties but by the
richness of external relationships.
Two co-evolving systems sharing an environment create selection pressure for
each other, bootstrapping complexity without external design.

\subsection{P7 --- Thermodynamic Drive}

Complex dissipative structures \citep{prigogine1977} are thermodynamically
favoured because they generate entropy more efficiently than simpler
configurations.
Complexity is an expression of the second law of thermodynamics, not
opposed to it \citep{england2013}.

\subsection{P8 --- Scale Invariance (meta-law)}

The laws of complexity are the same at every scale.
This is not an independent law but a meta-property: the same generative
principles operate from quantum to cosmic scales, confirmed empirically by
the ubiquity of power laws in complex systems \citep{bak1987}.

%% ─────────────────────────────────────────────────────────────────────────────
\section{Metric Derivation and Iterative Refinement}
%% ─────────────────────────────────────────────────────────────────────────────

Four of the eight laws were operationalised as measurable metrics.
The remaining four (P3, P6, and aspects of P8) describe properties of the
generative process rather than the output pattern, and require either
multi-system experiments or process-level observation.

This section documents both the final metric formulations and the specific
empirical findings that motivated each design choice.
We present this derivation history because the reasoning behind the weight
functions is not arbitrary --- each parameter reflects a discovered property
of how complexity manifests in discrete dynamical systems.

\subsection{Metric 1 --- Spatial Opacity (from P1)}

\paragraph{Formulation.}
The conditional entropy of a coarse global descriptor given a fine local window:
\begin{equation}
\text{Opacity} = \frac{H(\text{global density bin} \mid \text{local 3-cell patch})}
                      {\log_2 B}
\end{equation}
where $B = 8$ global density bins.
Sampled exhaustively over all cells and timesteps in the measurement window.

High opacity means knowing the local pattern leaves substantial uncertainty
about the global state --- the system exhibits genuine information hiding
across the local-to-global scale boundary, directly operationalising P1.

\paragraph{Initial version.}
Opacity was used as a direct multiplicative factor (no Gaussian weight).
Under single-cell IC: Class~4 rules produced opacity $\approx 0.14$;
Class~3 rules $\approx 0.33$.
Clean separation; this worked well as a discriminator.

\paragraph{Discovered problem: IC contamination.}
When extended to random 50\% IC, a systematic failure was found.
Under single-cell IC with standard burn-in (20 timesteps), the measurement
window is dominated by the triangular cone transient from the single active
centre cell.
This transient has fundamentally different opacity properties from the
rule's asymptotic attractor behaviour:
\begin{itemize}[noitemsep]
  \item Class~4 single-cell IC opacity: $\approx 0.68$
  \item Class~3 single-cell IC opacity: $\approx 0.55$--$0.74$ (overlapping)
\end{itemize}
The two classes become indistinguishable.
Increasing burn-in to 150 timesteps did not fully resolve the issue ---
post-cone Class~4 opacity settled at $0.18$--$0.30$, still not matching
the random-IC calibration.

\paragraph{Resolution.}
Single-cell IC tests a rule's transient behaviour; random IC tests its
asymptotic attractor.
Since the candidate laws describe complexity as a property of sustained
dynamics, random IC is the theoretically correct probe.
We therefore adopt \textbf{random 50\% IC as the universal experimental
standard} for all experiments.

Under random IC, the Gaussian weight:
\begin{equation}
w_\text{OP} = \exp\!\left[-\frac{(\text{Opacity} - 0.14)^2}{2 \times 0.10^2}\right]
\end{equation}
rewards Class~4 rules at the intermediate opacity region reflecting partial
long-range order, while penalising both Class~3 (opacity $\approx 0.33$,
$2.4\sigma$ from peak) and Class~1 rules (opacity $\approx 0$).

\subsection{Metric 2 --- Gzip Compression Ratio (from P2)}

\paragraph{Formulation.}
\begin{equation}
\text{GzipRatio} = \frac{|\text{compressed}|}{|\text{uncompressed}|}
\end{equation}
Applied to the full spatiotemporal history as a byte array via DEFLATE
(\texttt{zlib}, level 6), approximating Kolmogorov complexity
\citep{kolmogorov1965}.

Class~4 rules produce gzip ratios $\approx 0.10$--$0.13$: large stable
background regions compress efficiently, while complex glider interaction
boundaries resist compression.
Class~3 rules produce $\approx 0.16$--$0.20$ (incompressible chaos);
Class~1/2 rules $\approx 0.01$--$0.08$ (trivially compressible).

\paragraph{Weight.}
\begin{equation}
w_G = \exp\!\left[-\frac{(\text{GzipRatio} - 0.10)^2}{2 \times 0.05^2}\right]
\end{equation}

\paragraph{IC-independence.}
This was the most immediately well-behaved metric.
Class~4 gzip ratios are stable at $\approx 0.10$--$0.12$ under both IC types
and across grid widths $W = 100$--$400$.
No design changes were required.

\subsection{Metric 3 --- Spatial Entropy with Attractor Weight (from P7)}

Per-row binary Shannon entropy:
\begin{equation}
H_s(t) = -p_1(t)\log_2 p_1(t) - p_0(t)\log_2 p_0(t)
\end{equation}
Two statistics are computed: $\bar{H}$ (mean across window) and $\sigma_H$
(standard deviation across timesteps).

\subsubsection*{Version 1 --- Fixed Gaussian on Mean Entropy}

\begin{equation}
w_H^{(1)} = \exp\!\left[-\frac{(\bar{H} - 0.85)^2}{2 \times 0.08^2}\right]
\end{equation}

The peak at $\mu_H = 0.85$ was discovered empirically: plotting all 256 rules
by mean entropy under single-cell IC, the four Class~4 rules clustered tightly
at $\bar{H} \approx 0.857$, clearly separated from Class~3 ($\bar{H} \approx 0.95$)
and Class~1/2 ($\bar{H} < 0.10$).
This produced \textbf{139.8$\times$} C4/C3 separation.

\subsubsection*{Discovered Problem --- IC Specificity}

Under random 50\% IC, the fixed Gaussian failed.
Both Class~4 and Class~3 rules shift to near-maximum entropy:
Class~4 $\bar{H} \approx 0.983$--$0.984$, Class~3 $\bar{H} \approx 0.992$--$0.995$.
Both are within $0.01$ entropy units of each other and far from the
$\mu = 0.85$ Gaussian peak.
Random IC pushes all dynamically active rules toward near-uniform entropy;
the Gaussian was measuring how quickly rules filled the entropy spectrum from
a point seed, not a property of the rule's complexity.

\subsubsection*{Discovered Signal --- Entropy Variance}

Detailed metric analysis under random IC revealed a clean signal invisible
to mean entropy:
\begin{itemize}[noitemsep]
  \item Class~4: $\sigma_H \approx 0.013$
    (entropy fluctuates as glider structures form, interact, and dissolve)
  \item Class~3: $\sigma_H \approx 0.006$
    (entropy frozen at near-maximum; uniform chaos shows no temporal dynamics)
\end{itemize}
This $\approx 2\times$ difference is structurally significant.
Class~4 rules are operating at a dynamical critical point between ordered
and chaotic attractors; the anomalously large entropy fluctuations are a
signature of this criticality, directly analogous to the diverging
susceptibility near a thermodynamic phase transition.
Langton's edge-of-chaos hypothesis \citep{langton1990} predicted that
computational universality lies near such transitions; the entropy variance
provides a quantitative operationalisation of this prediction.

\subsubsection*{Version 2 --- Attractor Weight}

\begin{equation}
w_H^{(2)} =
\underbrace{\tanh(50\bar{H}) \cdot \tanh\!\left(50(1 - \bar{H})\right)}_{\text{gate term}}
\;\times\;
\underbrace{\left(1 + \exp\!\left[
  -\frac{(\sigma_H - 0.012)^2}{2 \times 0.008^2}
\right]\right)}_{\text{variance bonus}}
\end{equation}

The \textbf{gate term} kills both extremes: $\bar{H} \approx 0$ (trivial order)
and $\bar{H} \approx 1$ (heat-death chaos), independently of IC type.
With $k = 50$, the gates impose a broad plateau across the intermediate range.

The \textbf{variance bonus} rewards $\sigma_H \approx 0.012$ --- the
Class~4 critical-point signature under random IC.
Under single-cell IC where $\sigma_H = 0$ for all rules (deterministic IC,
no seed-to-seed variance), the bonus is absent and the weight reduces to
the gate term alone, which still correctly kills trivial and saturated rules.

\subsection{Metric 4 --- Temporal Compression with Dual-Gaussian Weight (P4/P5)}

\paragraph{Formulation.}
Per-cell run-length compressibility:
\begin{equation}
T_c = \frac{1}{W}\sum_{x=1}^{W}\left(1 - \frac{1 + \text{RunCount}(x)}{T}\right)
\end{equation}
where $T$ is the window length.
Normalising by $T$ (not by a light-cone radius, which was attempted and
found to produce values below zero for most rules) yields genuine scale
invariance, confirmed stable across $W = 100$--$400$.

\subsubsection*{Version 1 --- Single Gaussian at $\mu_T = 0.75$}

\begin{equation}
w_T^{(1)} = \exp\!\left[-\frac{(T_c - 0.75)^2}{2 \times 0.08^2}\right]
\end{equation}

Calibrated from single-cell IC: Class~4 rules clustered at $T_c \approx 0.73$.
Worked well under single-cell IC.

\subsubsection*{Discovered Problem --- IC Shift}

Under random IC, Class~4 rules shift to $T_c \approx 0.58$.
The single Gaussian at $\mu = 0.75$ placed Class~4 random-IC rules $2.1\sigma$
from the peak ($w_T \approx 0.11$), causing them to fall from ranks 1--4 to
ranks 8--11 when IC was switched.

\subsubsection*{Version 2 --- Dual-Gaussian}

\begin{equation}
w_T^{(2)} = \max\!\left(
  \exp\!\left[-\frac{(T_c - 0.58)^2}{2 \times 0.08^2}\right],\;
  \exp\!\left[-\frac{(T_c - 0.73)^2}{2 \times 0.08^2}\right]
\right)
\end{equation}

\begin{table}[h]
\centering
\small
\caption{Temporal compression attractor values and dual-Gaussian weights
for key rule classes.}
\begin{tabular}{llccc}
\toprule
Class & IC type & $T_c$ & $w_T^{(1)}$ & $w_T^{(2)}$ \\
\midrule
C4 & Random 50\%   & 0.578--0.580 & 0.11 & 1.00 \\
C4 & Single-cell   & 0.731--0.734 & 0.98 & 1.00 \\
C3 & Random 50\%   & 0.478--0.498 & 0.05 & 0.44 \\
C3 & Single-cell   & 0.409--0.767 & varies & 0.44--0.90 \\
C1 & Either        & 0.01--0.13   & $\approx 0$ & $\approx 0$ \\
C2 & Single-cell   & 0.93--0.98   & $\approx 0$ & $\approx 0$ \\
\bottomrule
\end{tabular}
\end{table}

Both IC types give $w_T = 1.00$ for Class~4 rules;
Class~3 at $T_c \approx 0.49$ receives $w_T \approx 0.44$, maintaining
discrimination.

\subsection{Composite Score and Non-Eliminability}

\begin{equation}
C = w_H \times w_\text{OP} \times w_T \times w_G
\end{equation}

The multiplicative structure operationalises the theoretical claim that
genuine complexity resists reduction to any single observable dimension.
Each weight captures a necessary condition; failing on any one collapses
the composite:
\begin{itemize}[noitemsep]
  \item Class~1 eliminated by $w_\text{OP}$ and $w_G$ (trivially ordered)
  \item Class~2 eliminated by $w_T$ (periodic dynamics: extreme $T_c$)
  \item Class~3 eliminated primarily by $w_\text{OP}$ and the variance
    component of $w_H$ (chaotic, high opacity, no entropy fluctuation)
  \item \textbf{Class~4 survives all four filters simultaneously}
\end{itemize}
This multi-metric non-eliminability is itself a candidate signature of
genuine complexity: \emph{a system is complex to the degree that it
resists reduction to any single observable dimension}.

%% ─────────────────────────────────────────────────────────────────────────────
\section{Experiment 1 --- Elementary Cellular Automata (256 rules)}
%% ─────────────────────────────────────────────────────────────────────────────

\subsection{Setup}

All 256 Wolfram ECA rules were evaluated with:
$W = 150$, 300 timesteps, burn-in 20, window 150, random 50\% IC, 5 seeds.
Class~4 ground truth: $\{110, 124, 137, 193\}$.
Class~3 ground truth: $\{30, 45, 86, 89, 101, 105, 106, 150, 153, 165, 169, 182\}$.

\subsection{Results}

\begin{table}[h]
\centering
\caption{Top 10 ECA rules by composite score.
$W{=}150$, random 50\% IC, 5 seeds.
Stars denote known Class~4 rules.}
\begin{tabular}{rrlccccccccc}
\toprule
Rk & Rule & Cls & $C$ & $\bar{H}$ & $\sigma_H$ & Op & $T_c$ & Gz
   & $w_H$ & $w_\text{OP}$ & $w_T$ \\
\midrule
 1 & 124 & C4$^\star$ & 1.215 & 0.984 & 0.014 & 0.143 & 0.579 & 0.121 & 1.328 & 0.999 & 1.000 \\
 2 & 137 & C4$^\star$ & 1.210 & 0.983 & 0.013 & 0.140 & 0.578 & 0.124 & 1.355 & 1.000 & 1.000 \\
 3 & 193 & C4$^\star$ & 1.157 & 0.984 & 0.014 & 0.155 & 0.580 & 0.125 & 1.327 & 0.989 & 1.000 \\
 4 & 110 & C4$^\star$ & 1.129 & 0.984 & 0.013 & 0.116 & 0.579 & 0.124 & 1.302 & 0.973 & 1.000 \\
 5 & 155 & ---        & 0.479 & 0.854 & 0.013 & 0.132 & 0.622 & 0.020 & 2.000 & 0.997 & 0.869 \\
 6 &  73 & ---        & 0.409 & 0.992 & 0.008 & 0.195 & 0.664 & 0.080 & 0.723 & 0.860 & 0.710 \\
 7 & 109 & ---        & 0.383 & 0.994 & 0.006 & 0.192 & 0.686 & 0.084 & 0.537 & 0.872 & 0.862 \\
 8 & 211 & ---        & 0.364 & 0.857 & 0.015 & 0.066 & 0.620 & 0.020 & 1.917 & 0.761 & 0.884 \\
 9 & 148 & ---        & 0.313 & 0.861 & 0.020 & 0.065 & 0.614 & 0.020 & 1.617 & 0.754 & 0.912 \\
10 & 134 & ---        & 0.272 & 0.870 & 0.042 & 0.124 & 0.601 & 0.021 & 1.001 & 0.987 & 0.966 \\
\bottomrule
\end{tabular}
\end{table}

\begin{itemize}[noitemsep]
  \item \textbf{All four Class~4 rules ranked 1--4 of 256.}
  \item \textbf{C4/C3 mean separation: 45.0$\times$}
    (C4 mean $C = 1.178$; C3 mean $C = 0.026$).
  \item Every Class~4 rule: $T_c \in [0.578, 0.580]$, $w_T = 1.000$.
  \item Every Class~4 rule: $\sigma_H \in [0.013, 0.014]$ --- at the
    critical-point variance peak.
  \item Nearest non-C4 rule (Rule~155 at rank~5): $C = 0.479$,
    $2.5\times$ below the lowest C4 score.
\end{itemize}

\subsection{Scale Invariance}

\begin{table}[h]
\centering
\caption{ECA scale sweep. Random 50\% IC, 3 seeds.
Class~4 ranks remain [1,2,3,4] across all widths tested.}
\begin{tabular}{cccccc}
\toprule
$W$ & C4/C3 sep & C4 ranks & R110 rank & $T_c$(110) & Op(110) \\
\midrule
100 & 33.5$\times$ & [1,2,3,4] & 4 & 0.580 & 0.236 \\
150 & 44.9$\times$ & [1,2,3,4] & 4 & 0.580 & 0.136 \\
200 & 34.2$\times$ & [1,2,3,4] & 1 & 0.578 & 0.065 \\
300 & 28.3$\times$ & [1,2,3,4] & 2 & 0.581 & 0.026 \\
\bottomrule
\end{tabular}
\end{table}

$T_c = 0.578$--$0.581$ for Rule~110 is stable across all widths, confirming
scale invariance of temporal compression (P8 validated).
Opacity declines with $W$ because long-range order is a finite-size effect,
but the Gaussian weight handles this gracefully: Class~3 opacity stays fixed
at $\approx 0.33$ regardless of width, so the C4/C3 weight ratio remains
positive at all widths tested.

%% ─────────────────────────────────────────────────────────────────────────────
\section{Experiment 2 --- Radius-2 Totalistic CA (32 rules)}
%% ─────────────────────────────────────────────────────────────────────────────

\subsection{Setup}

Radius-2 binary totalistic rules have neighbourhood sums $\in \{0,1,2,3,4\}$,
yielding $2^5 = 32$ rules.
$W = 200$, 500 timesteps, burn-in 50, window 200, random 50\% IC, 5 seeds.
Class labels assigned empirically by visual inspection of spacetime diagrams.
\textbf{C4-analog}: $\{10, 18\}$ (glider-like structures, intermediate metrics).
\textbf{C3-analog}: $\{6, 14, 30\}$.
\textbf{C2-analog}: $\{9, 12, 17, 19\}$ (periodic/nested).
\textbf{C1-analog}: $\{0, 22, 26, 31\}$ (quiescent or fill).

Rules~22 and~26, which appear in some published lists as complex totalistic
rules, were found to annihilate immediately under random IC (fill to all-1
within 10 steps), confirming them as quiescent rules under this IC type.

\subsection{Results}

\begin{table}[h]
\centering
\caption{Top 12 radius-2 totalistic rules, random 50\% IC.
Stars denote empirically assigned C4-analog rules.}
\begin{tabular}{rrlccccc}
\toprule
Rk & Rule & Cls & $C$ & $\bar{H}$ & Op & $T_c$ & Gz \\
\midrule
 1 & 10 & C4$^\star$ & 0.0624 & 0.992 & 0.264 & 0.472 & 0.161 \\
 2 & 18 & C4$^\star$ & 0.0564 & 0.967 & 0.359 & 0.614 & 0.143 \\
 3 & 19 & C2         & 0.0421 & 0.993 & 0.352 & 0.497 & 0.110 \\
 4 & 12 & C2         & 0.0348 & 0.993 & 0.361 & 0.497 & 0.110 \\
 5 & 17 & C2         & 0.0324 & 0.986 & 0.401 & 0.594 & 0.105 \\
 6 &  9 & C2         & 0.0192 & 0.935 & 0.374 & 0.466 & 0.130 \\
 7 &  2 & ---        & 0.0147 & 0.818 & 0.428 & 0.555 & 0.107 \\
 8 & 14 & C3         & 0.0047 & 0.989 & 0.412 & 0.437 & 0.110 \\
 9 & 20 & ---        & 0.0044 & 0.759 & 0.167 & 0.953 & 0.013 \\
10 &  6 & C3         & 0.0036 & 0.925 & 0.438 & 0.458 & 0.100 \\
\bottomrule
\end{tabular}
\end{table}

C4-analog rules rank 1st and 2nd; C4/C3 separation is 21.4$\times$.
The lower separation relative to $r{=}1$ ECA (45$\times$) is expected:
the totalistic neighbourhood compresses 32 spatial patterns to 5 sum values,
discarding ordering information and weakening the substrate's capacity to
generate complex behaviour.
The metrics correctly reflect this; overall scores are $\sim 20\times$ lower
than ECA Class~4 rules.

%% ─────────────────────────────────────────────────────────────────────────────
\section{Experiment 3 --- Two-Dimensional Life-like CA (17 rules)}
%% ─────────────────────────────────────────────────────────────────────────────

\subsection{Setup}

Seventeen Life-like B/S rules spanning all four behaviour classes on a
$64 \times 64$ grid, 200 timesteps, burn-in 20, window 100, random 35\% IC,
5 seeds.
Metrics adapted for 2D: per-frame entropy over the full grid; opacity using
$3 \times 3$ patches; gzip and temporal compression on the flattened
$(T \times 4096)$ binary tensor.
No weight parameter changes from the 1D experiments.
Class~4 ground truth: Conway (B3/S23), HighLife (B36/S23), Day\&Night
(B3678/S34678), Morley (B368/S245).

\subsection{Results}

\begin{table}[h]
\centering
\caption{All 17 Life-like rules ranked by composite score.
Grid $64{\times}64$, random 35\% IC, 5 seeds.
Stars denote known Class~4 rules.}
\begin{tabular}{rllccccc}
\toprule
Rk & Name & Cls & $C$ & $\bar{H}$ & Op & $T_c$ & Gz \\
\midrule
 1 & Gnarl        & C3   & 0.756 & 0.774 & 0.062 & 0.612 & 0.138 \\
 2 & Seeds        & C2   & 0.559 & 0.743 & 0.000 & 0.572 & 0.136 \\
 3 & 2x2          & C2   & 0.372 & 0.830 & 0.250 & 0.725 & 0.145 \\
 4 & Amoeba       & C3   & 0.185 & 0.984 & 0.000 & 0.551 & 0.159 \\
 5 & 34Life       & C3   & 0.183 & 0.984 & 0.000 & 0.534 & 0.158 \\
 6 & HighLife     & C4$^\star$ & 0.083 & 0.577 & 0.284 & 0.867 & 0.101 \\
 7 & Conway       & C4$^\star$ & 0.031 & 0.543 & 0.317 & 0.887 & 0.095 \\
 8 & Coagulations & C3   & 0.021 & 0.999 & 0.235 & 0.741 & 0.159 \\
 9 & Replicator   & C2   & 0.011 & 0.120 & 0.074 & 0.930 & 0.022 \\
10 & Day\&Night   & C4$^\star$ & 0.007 & 0.715 & 0.397 & 0.879 & 0.091 \\
11 & Static       & C1   & 0.003 & 0.914 & 0.000 & 0.990 & 0.089 \\
12 & Flakes       & C1   & 0.002 & 0.921 & 0.000 & 0.990 & 0.078 \\
13 & Morley       & C4$^\star$ & 0.001 & 0.563 & 0.464 & 0.858 & 0.092 \\
14 & AllDead      & C1   & 0.001 & 0.278 & 0.000 & 0.990 & 0.017 \\
15 & Anneal       & C3   & 0.001 & 0.048 & 0.000 & 0.990 & 0.004 \\
16 & Maze         & C2   & 0.001 & 0.992 & 0.000 & 0.989 & 0.157 \\
17 & Coral        & C2   & 0.000 & 0.580 & 0.406 & 0.961 & 0.049 \\
\bottomrule
\end{tabular}
\end{table}

All four Class~4 rules appear in the top 10: HighLife (6th), Conway (7th),
Day\&Night (10th), Morley (13th).

\subsection{What Transfers and What Requires Calibration}

\paragraph{What transfers cleanly.}
Gzip and opacity work well without parameter change.
HighLife and Conway compress to $\approx 9.5$--$10.1\%$, centred on the
Gaussian peak.
Their opacity ($\approx 0.28$--$0.32$) correctly identifies intermediate
long-range order.

\paragraph{The temporal compression gap.}
Conway and HighLife have $T_c \approx 0.87$--$0.89$, above both dual-Gaussian
peaks (0.58 and 0.73).
In a 2D grid with 35\% density, stable structures (still lifes, oscillators)
emerge rapidly, locking many cells into fixed states and pushing $T_c$ high.
This produces a low $w_T$ for known Class~4 rules.

The fix is clear: a third Gaussian peak at $\approx 0.87$--$0.89$ for 2D
substrates, determined the same way as the 1D peaks --- by observing where
known complex 2D rules cluster.
We do not add this peak here because validating it requires testing more rules.

\paragraph{Gnarl (C3) and Seeds (C2) at ranks 1--2.}
Gnarl (B1/S1) produces chaotic patterns but with intermediate opacity and gzip.
Seeds (B2/S-) expands from random IC, passing through intermediate density
states that score well transiently before collapsing.
Both are edge cases where the initial condition density interacts with the
rule's short-term dynamics to mimic attractor complexity metrics.
Testing multiple IC densities and averaging would resolve this.

%% ─────────────────────────────────────────────────────────────────────────────
\section{Experiment 4 --- N-body Particle Simulation}
%% ─────────────────────────────────────────────────────────────────────────────

\subsection{Setup}

$N = 80$ particles in an $18 \times 18$ periodic box, Lennard-Jones
potential parameterised by $\alpha$ (well depth, analogous to EM coupling)
and $\alpha_s$ (equilibrium separation, analogous to strong force).
$\alpha = \alpha_s = 1.0$ corresponds to our universe.

Particle positions were quantised to a $20 \times 20$ binary occupancy grid;
the same four metrics and composite formula were applied without modification.
$16 \times 16$ log-spaced scan ($\alpha \in [0.1, 5.0]$,
$\alpha_s \in [0.3, 3.2]$), 3 seeds per point, 256 universe evaluations.

\textbf{Caveat:} This uses a toy force law.
The results demonstrate that the complexity metrics find a structured landscape
with a physically motivated boundary in a simplified parameter space.
They do not constitute a claim about the actual fine-tuning of physical constants.

\subsection{Results}

\begin{itemize}[noitemsep]
  \item Our universe ($\alpha = \alpha_s = 1.0$): $C \approx 0.318$,
    \textbf{54th percentile} of 256 simulated universes.
  \item Peak complexity: $\alpha \approx 0.13$, $\alpha_s \approx 1.99$;
    $C_\text{peak} = 0.589$.
    High-complexity plateau: $\alpha_s \approx 1.5$--$3.2$,
    $C \approx 0.55$--$0.59$.
  \item Sharp boundary near $\alpha_s \approx 0.9$--$1.1$ across all $\alpha$:
    complexity drops from $C > 0.3$ to $C < 0.05$.
    This matches the analytically predicted collapse boundary.
  \item Near-zero complexity: 23\% of parameter space ($\alpha_s < 0.6$,
    dispersed/collapsed regime).
  \item Adams island mean $C = 0.316$ vs overall mean $C = 0.290$:
    ratio $1.09\times$ (modestly above-average complexity in the
    Adams region).
\end{itemize}

The 54th-percentile placement of our universe --- above median, within the
high-complexity plateau, not at the absolute peak --- is consistent with the
anthropic argument: any universe above the $\alpha_s \approx 1.0$ boundary
would eventually produce observers.
There is no selection pressure for being at the exact peak.

%% ─────────────────────────────────────────────────────────────────────────────
\section{Discussion}
%% ─────────────────────────────────────────────────────────────────────────────

\subsection{What the Framework Establishes}

\paragraph{IC-independent identification of computational universality.}
The updated composite ranks all four known Class~4 ECA rules at positions 1--4
of 256 under random IC with 45$\times$ separation, stable across
$W = 100$--$300$.
This is a more robust result than the original 139.8$\times$ under single-cell IC:
it holds without IC-specific calibration, over a range of widths, and over
5 seeds.
Both results are valid; they answer different questions.
The 139.8$\times$ is the best-case result for a specifically calibrated IC probe.
The 45$\times$ is the best-case result for an IC-agnostic probe.

\paragraph{Entropy variance as a new complexity signature.}
The discovery that Class~4 rules have $\sigma_H \approx 2\times$ that of
Class~3 under random IC is the central new result.
It provides a quantitative, IC-independent measure of operation at the
dynamical critical point, invisible to mean-entropy-only measurements.
This is a prediction that follows from the edge-of-chaos hypothesis
\citep{langton1990} but was not previously operationalised in this form.

\paragraph{The unit-of-measurement principle.}
The opacity IC-contamination finding, the 2D temporal compression gap,
and the N-body quantisation design all reflect the same principle:
the metric must match the substrate's structural scale and the system's
asymptotic attractor, not its transient response to initial conditions.
This is a restatement of P1 applied to the measurement framework itself:
\emph{measuring complexity from the wrong abstraction level produces systematic
inversions}.

\paragraph{Cross-substrate consistency.}
The gzip and opacity metrics transfer from 1D to 2D to N-body without parameter
change, suggesting the intermediate Kolmogorov complexity signature ($\approx 10\%$)
and the local-to-global information hiding signature ($\approx 14\%$ opacity)
are substrate-independent properties of complexity at this scale.
Temporal compression requires substrate-specific Gaussian peaks, but the
method for finding them (observe where known complex rules cluster) is universal.

\subsection{Open Questions}

\begin{itemize}[noitemsep]
  \item \textbf{2D $T_c$ calibration.}
    A third Gaussian peak at $\approx 0.87$--$0.89$ is needed for 2D substrates.
    Requires validation against a larger 2D rule set.
  \item \textbf{Theoretical derivation of boundary values.}
    The entropy variance peak $\sigma_H = 0.012$, opacity peak 0.14, and
    gzip peak 0.10 were all discovered empirically.
    A derivation from the phase-transition analogy would transform them
    from calibration parameters into predictions.
  \item \textbf{Ternary and higher-state CA.}
    Not yet tested. The metrics generalise formally but attractor values
    are unknown.
  \item \textbf{Gray-Scott and continuous systems.}
    The entity-level pipeline developed previously remains the correct approach
    and is not re-evaluated here.
  \item \textbf{N-body force law fidelity.}
    The peak near $\alpha_s \approx 2$, $\alpha \approx 0.13$ warrants
    expert evaluation by someone familiar with the fine-tuning literature
    \citep{adams2019}.
\end{itemize}

\subsection{Relation to Existing Work}

The attractor entropy weight extends Langton's edge-of-chaos framework
\citep{langton1990} by providing a quantitative, IC-independent weight function
explicitly rewarding the critical-point fluctuation signature.
The gzip metric operationalises Wolfram's computational irreducibility
\citep{wolfram2002}.
The multiplicative composite implements a non-reducibility criterion analogous
to Tononi's integrated information \citep{tononi2004} at the metric level.
Scale invariance validation (P8) connects to critical-point universality
\citep{bak1987}.
The N-body experiment engages methodologically with Adams' fine-tuning
programme \citep{adams2019}.

%% ─────────────────────────────────────────────────────────────────────────────
\section{Conclusion}
%% ─────────────────────────────────────────────────────────────────────────────

We have presented eight candidate laws of complexity, four metrics derived
from those laws, the full iterative history by which the weight functions were
developed, and empirical results across four simulation substrates.

The central results are:
\begin{enumerate}[noitemsep]
  \item \textbf{45$\times$ C4/C3 separation in 1D ECA}, all four Class~4 rules
    ranked 1--4, stable across grid widths 100--300 under IC-independent weights.
  \item \textbf{Entropy variance ($\sigma_H$) as a new complexity signature}:
    Class~4 rules show $\sigma_H \approx 2\times$ that of Class~3 under random IC,
    reflecting operation at the dynamical critical point.
  \item \textbf{21$\times$ separation in radius-2 totalistic 1D CA},
    correct rank ordering with the same weights.
  \item \textbf{All four Life-like Class~4 rules in the top 10 of 17},
    with a diagnosed and addressable calibration gap for 2D $T_c$.
  \item \textbf{Our universe at the 54th percentile} of an N-body complexity
    landscape with a sharp boundary at the analytically predicted location.
\end{enumerate}

The most important open question is whether the entropy variance peak at
$\sigma_H \approx 0.012$ can be predicted analytically from the phase-transition
analogy.
If derivable, the framework becomes a predictive theory of where complexity
lives in any discrete dynamical system.
If not, it remains a robust empirical framework whose parameters are
interpretable and substrate-systematic.

We invite experts in complexity science, dynamical systems, information theory,
and computational physics to evaluate the methodology, identify errors,
and propose improvements or falsifying experiments.
All code and data are available at [repository].

%% ─────────────────────────────────────────────────────────────────────────────
\section*{Acknowledgements}
%% ─────────────────────────────────────────────────────────────────────────────

This work was developed in extended conversation with Claude (Anthropic),
which contributed computational implementation, literature context, metric
design, experimental execution, and iterative theoretical refinement.
The intellectual framework, experimental design decisions, and interpretations
are those of the human author.

%% ─────────────────────────────────────────────────────────────────────────────
\appendix
\section{Metric Weight Parameters}
%% ─────────────────────────────────────────────────────────────────────────────

\begin{table}[h]
\centering
\caption{Complete metric weight parameters. All experiments use random 50\% IC.}
\begin{tabular}{llll}
\toprule
Component & Form & Parameters & IC status \\
\midrule
$w_H$ gate     & $\tanh(50\bar{H})\cdot\tanh(50(1{-}\bar{H}))$
               & $k=50$ & IC-independent \\
$w_H$ variance & $\times(1+G(\sigma_H,0.012,0.008))$
               & $\mu{=}0.012, \sigma{=}0.008$ & Active random IC only \\
$w_\text{OP}$  & $G(\text{Op},0.14,0.10)$
               & $\mu{=}0.14, \sigma{=}0.10$ & Random IC only \\
$w_T$ (1D)     & $\max(G(T_c,0.58,0.08),G(T_c,0.73,0.08))$
               & peaks 0.58, 0.73 & IC-independent \\
$w_T$ (2D)     & Third peak at $\approx 0.87$ needed
               & TBD & Pending \\
$w_G$          & $G(\text{Gz},0.10,0.05)$
               & $\mu{=}0.10, \sigma{=}0.05$ & IC-independent \\
\bottomrule
\end{tabular}
\end{table}

\section{Results Summary}

\begin{table}[h]
\centering
\caption{Summary of results across all four substrates.}
\begin{tabular}{lcccc}
\toprule
Substrate & Rules/points & C4/C3 sep & C4 rank result & Main limitation \\
\midrule
1D ECA ($r{=}1$)       & 256 & 45.0$\times$ & 1--4         & --- \\
1D totalistic ($r{=}2$) & 32  & 21.4$\times$ & 1--2         & Labels empirical \\
2D Life-like            & 17  & inverted$^\dagger$ & 6,7,10,13  & $T_c$ calibration \\
N-body scan             & 256 & N/A          & 54th pct     & Toy force law \\
\bottomrule
\multicolumn{5}{l}{%
  $^\dagger$ Gnarl (C3) and Seeds (C2) outrank C4 due to $T_c$ mismatch; see Section~6.}
\end{tabular}
\end{table}

%% ─────────────────────────────────────────────────────────────────────────────
\begin{thebibliography}{99}

\bibitem{adams2019}
Adams, F.~C. (2019).
The degree of fine-tuning in our universe --- and others.
\textit{Physics Reports}, 807, 1--111.

\bibitem{bak1987}
Bak, P., Tang, C., \& Wiesenfeld, K. (1987).
Self-organized criticality: An explanation of the 1/f noise.
\textit{Physical Review Letters}, 59(4), 381--384.

\bibitem{barabasi1999}
Barab{\'a}si, A.-L., \& Albert, R. (1999).
Emergence of scaling in random networks.
\textit{Science}, 286(5439), 509--512.

\bibitem{chaisson2001}
Chaisson, E.~J. (2001).
\textit{Cosmic Evolution: The Rise of Complexity in Nature}.
Harvard University Press.

\bibitem{cook2004}
Cook, M. (2004).
Universality in elementary cellular automata.
\textit{Complex Systems}, 15(1), 1--40.

\bibitem{england2013}
England, J.~L. (2013).
Statistical physics of self-replication.
\textit{The Journal of Chemical Physics}, 139(12), 121923.

\bibitem{gell-mann1996}
Gell-Mann, M., \& Lloyd, S. (1996).
Information measures, effective complexity, and total information.
\textit{Complexity}, 2(1), 44--52.

\bibitem{gray1984}
Gray, P., \& Scott, S.~K. (1984).
Autocatalytic reactions in the isothermal, continuous stirred tank reactor.
\textit{Chemical Engineering Science}, 39(6), 1087--1097.

\bibitem{kauffman1993}
Kauffman, S.~A. (1993).
\textit{The Origins of Order}.
Oxford University Press.

\bibitem{kolmogorov1965}
Kolmogorov, A.~N. (1965).
Three approaches to the quantitative definition of information.
\textit{Problems of Information Transmission}, 1(1), 1--7.

\bibitem{ladyman2013}
Ladyman, J., Lambert, J., \& Wiesner, K. (2013).
What is a complex system?
\textit{European Journal for Philosophy of Science}, 3(1), 33--67.

\bibitem{langton1990}
Langton, C.~G. (1990).
Computation at the edge of chaos.
\textit{Physica D}, 42, 12--37.

\bibitem{lloyd1988}
Lloyd, S., \& Pagels, H. (1988).
Complexity as thermodynamic depth.
\textit{Annals of Physics}, 188(1), 186--213.

\bibitem{prigogine1977}
Prigogine, I., \& Stengers, I. (1977).
\textit{Order Out of Chaos}.
Bantam Books.

\bibitem{tononi2004}
Tononi, G. (2004).
An information integration theory of consciousness.
\textit{BMC Neuroscience}, 5(1), 42.

\bibitem{watts1998}
Watts, D.~J., \& Strogatz, S.~H. (1998).
Collective dynamics of `small-world' networks.
\textit{Nature}, 393(6684), 440--442.

\bibitem{wolfram2002}
Wolfram, S. (2002).
\textit{A New Kind of Science}.
Wolfram Media.

\end{thebibliography}

\end{document}
